\documentclass[11pt]{article}
\usepackage[margin=1in]{geometry}
\usepackage[T1]{fontenc}
\usepackage{mathptmx}\usepackage[scaled=.92]{helvet}\usepackage{courier}\usepackage{microtype}\renewcommand{\familydefault}{\sfdefault}
\usepackage{booktabs,tabularx,array,xcolor,graphicx,enumitem,fancyhdr,titlesec,listings,amsmath}
\usepackage[hidelinks]{hyperref}
\definecolor{accent}{HTML}{1F4E79}
\definecolor{soft}{HTML}{EEF3F8}
\definecolor{good}{HTML}{2E7D32}
\titleformat{\section}{\large\bfseries\color{accent}}{\thesection}{0.6em}{}
\titleformat{\subsection}{\normalsize\bfseries\color{accent}}{\thesubsection}{0.6em}{}
\setlength{\parindent}{0pt}\setlength{\parskip}{0.55em}
\renewcommand{\arraystretch}{1.25}
\pagestyle{fancy}\fancyhf{}
\fancyfoot[C]{\small\thepage}
\renewcommand{\headrulewidth}{0pt}
\lstset{basicstyle=\ttfamily\small,backgroundcolor=\color{soft},frame=none,breaklines=true,columns=fullflexible,keepspaces=true,xleftmargin=0.5em,aboveskip=0.6em,belowskip=0.4em}
\newcommand{\code}[1]{\texttt{\small\detokenize{#1}}}
\usepackage[most]{tcolorbox}
\newtcolorbox{callout}[1]{colback=soft,colframe=accent,boxrule=0.5pt,arc=2pt,left=6pt,right=6pt,top=4pt,bottom=4pt,title={#1},fonttitle=\bfseries}
\usepackage{colortbl}
\newcommand{\titleblock}[2]{\begin{center}{\LARGE\bfseries\color{accent}#1}\\[0.4em]{\large #2}\\[0.3em]{\small Kevin Williams \quad|\quad October 23, 2026}\end{center}\vspace{0.3em}}
\begin{document}
\titleblock{Volatility Surface Tool --- User Guide}{Type one ticker, get an interactive implied-volatility surface}

\section{What this tool does}
\code{vol_surface("AAPL")} downloads the current option chain for a stock, works out the \emph{implied volatility} of every liquid option, and draws three linked views: a 3-D surface, the ``smile'' for one expiry, and the term structure. A second window explains how to read all of it in plain language. Under the hood it connects a Python/OCaml option-pricing engine to MATLAB's symbolic and numerical tools.

\begin{center}\includegraphics[width=0.92\linewidth]{vol_surface_aapl.png}\end{center}

\section{Before you start}
\begin{tabularx}{\linewidth}{@{}lX@{}}
\toprule
\textbf{Need} & \textbf{Details}\\\midrule
MATLAB & R2020a or newer, with the Symbolic Math Toolbox and the Statistics and Machine Learning Toolbox\\
Python & A Python with \code{yfinance}, \code{pandas}, \code{numpy} and \code{scipy}. The tool looks in \code{<repo>/.venv}, then MATLAB's configured Python, then \code{python3}\\
OCaml engine & Optional. Adds a second, independent volatility for cross-checking. Turn off with \code{"UseOCaml", false}\\
Internet & Needed to fetch a chain from Yahoo Finance\\\bottomrule
\end{tabularx}

\section{Quick start}
\begin{enumerate}[leftmargin=1.5em]
\item In MATLAB, set the current folder to the repository's \code{matlab} folder.
\item Run one line:
\begin{lstlisting}
vol_surface("AAPL")
\end{lstlisting}
\item Wait 10--60 seconds for the download. Two windows open: the surface figure and \emph{How to read the surface}.
\end{enumerate}
Ticker symbols may contain letters, digits, ``.'' and ``-'' (up to 10 characters). The tool writes a dated snapshot to \code{snapshots/} so you can reuse it later.

\section{When to run it}
Yahoo Finance serves only the \emph{current} chain, and it is meaningful when the market is open (9:30--16:00 Eastern, weekdays). Before the open most bids and asks are zero.

\begin{tabularx}{\linewidth}{@{}lX@{}}
\toprule
\textbf{Situation} & \textbf{Command}\\\midrule
Market open (best) & \code{vol_surface("AAPL")}\\
Too early in the day & \code{vol_surface("AAPL", "Session", "prior")} rebuilds yesterday's close from each option's last trade\\
Not sure & \code{vol_surface("AAPL", "Session", "auto")} tries live, then falls back to the prior session\\
Offline or repeat a run & \code{vol_surface("AAPL", "Refresh", false)} reuses the saved snapshot\\\bottomrule
\end{tabularx}

\textbf{Prior-session mode} keeps only trades from the last \code{PriorWindowMin} minutes (default 60) of the previous session and prices them against that session's closing price. It also works shortly after the open: the tool measures the share of contracts that have already traded today and refuses if that exceeds 20\%, because liquid contracts trade first and the sample would skew toward illiquid strikes. Prior surfaces are noisier than live ones, and the figure title says so.

\section{Options}
\begin{tabularx}{\linewidth}{@{}>{\ttfamily\small}l>{\ttfamily\small}lX@{}}
\toprule
\textbf{Name} & \textbf{Default} & \textbf{Meaning}\\\midrule
Refresh & true & Fetch live data; false uses the cached snapshot. If a fetch fails and a cache exists, the cache is used with a warning\\
Session & "live" & "live", "prior" or "auto" (see above)\\
PriorWindowMin & 60 & Prior mode: keep trades from the last N minutes of that session\\
MaxExpiries & 8 & Number of expiries to download\\
MinDays & 7 & Skip expiries closer than this many days\\
Expiry & "" & Expiry shown in the smile panel, e.g. "2026-11-20". Default: the expiry with the most quotes\\
UseOCaml & true & Add the OCaml reference volatility and cross-check against it\\
Python & "" & Path to a specific Python. Default: automatic search\\
Explain & true & Open the plain-language explainer window\\\bottomrule
\end{tabularx}

Example combining several options:
\begin{lstlisting}
out = vol_surface("MSFT", "MaxExpiries", 6, "Expiry", "2026-11-20", "Explain", false);
\end{lstlisting}

\section{Reading the results}
\begin{itemize}[leftmargin=1.5em]
\item \textbf{Surface (left).} Height is implied volatility. The horizontal axes are how far the strike is from the current price (log-moneyness, 0 = at the forward) and days to expiry. Black dots are real quotes. Blank areas have no reliable data and are never extrapolated.
\item \textbf{Smile (top right).} One expiry cut out of the surface, with the market quotes as dots.
\item \textbf{Term structure (bottom right).} Volatility at the forward price across expiries.
\item \textbf{Explainer window.} Six panels that define options, volatility, implied volatility, the smile and the term structure using this ticker's numbers.
\end{itemize}
The optional return value \code{out} is a struct with \code{csv}, \code{chain}, \code{res}, \code{surface}, \code{smile}, \code{term}, \code{figure} and \code{explainFigure}, handy for your own analysis.

\section{Using the pieces separately}
Any chain CSV works, not only the tool's own snapshots.
\begin{lstlisting}
res = run_vol_checks("snapshots/AAPL_latest.csv");   % solve + compare to OCaml
s   = build_surface(res);                            % interpolated surface
sl  = smile_slice(s, res, "2026-11-20");             % one expiry
ts  = term_structure_slice(s, 0);                    % at-the-forward term structure
\end{lstlisting}
Required CSV columns (any case): \code{spot, expiry, T, r, type, strike} and either \code{mid} or both \code{bid} and \code{ask}. Optional: \code{q} (dividend yield, default 0), \code{symbol}, and \code{iv_ocaml} (reference volatility). \code{T} is years to expiry; \code{type} is call or put.

\section{Checking that it works}
\begin{lstlisting}
test_vol_pipeline     % pricing, solver and pipeline, no network needed
test_surface          % interpolation and slices against a known surface
verify_demo           % one-page pass/fail verification report
verify_demo("snapshots/AAPL_latest.csv")   % also compare a real snapshot to OCaml
\end{lstlisting}
The companion document \emph{Verification Process} explains each check.

\section{Troubleshooting}
\begin{tabularx}{\linewidth}{@{}>{\raggedright\arraybackslash}p{0.38\linewidth}X@{}}
\toprule
\textbf{Message or symptom} & \textbf{What to do}\\\midrule
``has only N usable out-of-the-money quote(s)'' & The chain is thin, usually before the open. Rerun when the market is open, use \code{"Session","prior"}, or pass \code{"Refresh", false}\\
``\% of contracts have already traded since the prior session closed'' & Too much of today's trading has overwritten yesterday's prices. Use live mode\\
``Python exporter failed'' & The Python found lacks \code{yfinance} or its relatives. Pass \code{"Python","/path/to/python"}\\
``No cached snapshot at \ldots'' & There is nothing saved for this ticker and mode. Run once with \code{"Refresh", true}\\
Dip in one short expiry & Few quotes with wide spreads; treat that expiry with caution\\
Unknown ticker or empty chain & Check the symbol; some tickers have no listed options\\\bottomrule
\end{tabularx}

\section{Limitations}
\begin{itemize}[leftmargin=1.5em]
\item Black-Scholes with a constant dividend yield treats American single-stock options as European; the tool uses only out-of-the-money quotes to limit the effect.
\item The surface is interpolated, not fitted with a no-arbitrage model. Treat it as a faithful picture of the quotes, not a tradable price list.
\item Quote data comes from Yahoo Finance and may be delayed or stale.
\item Educational use only; nothing here is investment advice.
\end{itemize}
\end{document}
